% =====================================================================
%  CARNET D'ENTRAÎNEMENT — Fiche 12 (Chimie organique)
%  THÈME 3 — Nomenclature et priorité des fonctions
%  NIVEAU FACILE — CORRIGÉ
% =====================================================================
\documentclass[11pt,a4paper]{article}
\usepackage[utf8]{inputenc}
\usepackage[T1]{fontenc}
\usepackage{lmodern}
\usepackage[french]{babel}
\usepackage{amsmath,amssymb,mathtools}
\providecommand{\degree}{^\circ}
\usepackage{icomma}
\usepackage{siunitx}
\sisetup{output-decimal-marker={,},per-mode=symbol}
\usepackage[version=4]{mhchem}
\usepackage{chemfig}
\setchemfig{atom sep=2em,bond length=0.9em,cram width=2.5pt}
\usepackage{xcolor}
\usepackage{array}
\usepackage{booktabs}
\usepackage{tabularx}
\usepackage[most]{tcolorbox}
\usepackage{enumitem}
\usepackage{tikz}
\usetikzlibrary{arrows.meta,positioning,calc}
\usepackage{fancyhdr}
\usepackage{titlesec}
\usepackage[hidelinks]{hyperref}
\usepackage[a4paper,margin=2.1cm,headheight=26pt,headsep=9pt]{geometry}

\definecolor{primary}{RGB}{150,98,162}
\definecolor{primarydark}{RGB}{92,52,110}
\definecolor{excol}{RGB}{0,135,90}
\definecolor{attncol}{RGB}{200,40,55}
\definecolor{methcol}{RGB}{90,90,100}

\newcommand{\runtitle}{Thème 3 — Facile — Corrigé}
\pagestyle{fancy}\fancyhf{}
\fancyhead[L]{\small\textcolor{primarydark}{\textbf{Fiche 12} — Chimie organique}}
\fancyhead[R]{\small\textcolor{primarydark}{\runtitle}}
\fancyfoot[C]{\small\thepage}
\renewcommand{\headrulewidth}{0.8pt}
\renewcommand{\headrule}{\hbox to\headwidth{\color{primary}\leaders\hrule height \headrulewidth\hfill}}
\setlength{\parindent}{0pt}\setlength{\parskip}{3pt}

\tcbset{boxbase/.style={enhanced,breakable,boxrule=0.9pt,arc=2.5pt,
   left=3mm,right=3mm,top=2mm,bottom=2mm,fonttitle=\bfseries,coltitle=white,
   toptitle=0.5mm,bottomtitle=0.5mm}}
\newtcolorbox[auto counter]{corr}[1]{boxbase,colback=excol!5,colframe=excol,
   title={\textsc{Corrigé}~\thetcbcounter\ \textmd{--- \itshape #1}}}
\newcommand{\rep}[1]{\textcolor{primarydark}{\textbf{#1}}}

\newcommand{\banner}[2]{%
\begin{tcolorbox}[enhanced,colback=excol,colframe=excol,arc=5pt,boxrule=0pt,
   left=5mm,right=5mm,top=2.5mm,bottom=2.5mm,coltext=white]
{\large\bfseries #1}\\[1pt]{\small #2}
\end{tcolorbox}\vspace{2pt}}

\begin{document}

\banner{Thème 3 — Nomenclature et priorité des fonctions}{Niveau \textbf{Facile} — Corrigé détaillé \quad$\vert$\quad Fiche 12, Chimie organique}

% ---------------------------------------------------------------------
\begin{corr}{donner le nom (fonction unique)}
\begin{itemize}[leftmargin=1.6em,itemsep=1pt]
  \item \textbf{(A)} $\ce{CH3CH2CH2OH}$ : alcool, 3 C, $\ce{-OH}$ en bout $\Rightarrow$ \rep{propan-1-ol} ;
  \item \textbf{(B)} $\ce{CH3CH2CH2CHO}$ : aldéhyde, 4 C (le $\ce{CHO}$ est le C1) $\Rightarrow$ \rep{butanal} ;
  \item \textbf{(C)} $\ce{CH3COCH3}$ : cétone, 3 C, $\ce{C=O}$ en position 2 $\Rightarrow$ \rep{propanone} (propan-2-one) ;
  \item \textbf{(D)} $\ce{CH3CH2COOH}$ : acide, 3 C (le $\ce{COOH}$ est le C1) $\Rightarrow$ \rep{acide propanoïque}.
\end{itemize}
\end{corr}

% ---------------------------------------------------------------------
\begin{corr}{numéroter la chaîne}
Pour un acide carboxylique, \textbf{le carbone de la fonction porte toujours le n\textsuperscript{o}1}. Donc :
\[ \underset{\text{n}^\circ 1}{\ce{COOH}}-\underset{2}{\ce{CH2}}-\underset{3}{\ce{CH2}}-\underset{4}{\ce{CH3}} \]
Le carbone du $\ce{-COOH}$ est le \rep{n\textsuperscript{o}1} ; on numérote ensuite $2,3,4$ jusqu'au $\ce{CH3}$ terminal. (C'est un cas imposé, valable aussi pour ester, aldéhyde et amide.)
\end{corr}

% ---------------------------------------------------------------------
\begin{corr}{repérer la fonction principale}
\textbf{a)} On a un $\ce{-CHO}$ (\rep{aldéhyde}) et un $\ce{-OH}$ (\rep{alcool}).\\
\textbf{b)} Dans l'échelle, $\ce{-CHO} > \ce{-OH}$ : la \rep{fonction principale est l'aldéhyde} (suffixe \textit{-al}). L'alcool, moins prioritaire, devient un \rep{préfixe « hydroxy- »}. (Le nom complet serait \emph{3-hydroxypropanal}.)
\end{corr}

% ---------------------------------------------------------------------
\begin{corr}{du suffixe à la fonction}
\begin{itemize}[leftmargin=1.6em,itemsep=1pt]
  \item butan-2-\textbf{ol} $\Rightarrow$ \rep{alcool} ;
  \item propan\textbf{al} $\Rightarrow$ \rep{aldéhyde} ;
  \item acide éthano\textbf{ïque} $\Rightarrow$ \rep{acide carboxylique} ;
  \item butan-2-\textbf{one} $\Rightarrow$ \rep{cétone} ;
  \item éthan\textbf{amide} $\Rightarrow$ \rep{amide}.
\end{itemize}
\end{corr}

% ---------------------------------------------------------------------
\begin{corr}{nommer une molécule ramifiée}
La fonction est un \textbf{alcool} (suffixe \textit{-ol}). La chaîne la plus longue contenant le $\ce{-OH}$ compte \textbf{3 carbones} (propan-). On numérote pour que le $\ce{-OH}$ ait le plus petit indice : $\ce{-OH}$ en \textbf{1}, la ramification $\ce{-CH3}$ tombe alors en \textbf{2}. Nom : \rep{2-méthylpropan-1-ol}.
\end{corr}

\end{document}
